% CCAI 2026 submission — double-blind, LNCS format, 8+2 pages
\documentclass[runningheads]{llncs}

\usepackage[T1]{fontenc}
\usepackage{graphicx}
\usepackage{amsmath,amssymb}
\usepackage{booktabs}
\usepackage{multirow}
\usepackage{xcolor}
\usepackage[textsize=tiny,disable]{todonotes}
\usepackage{comment}

% Hyperref last (standard practice)
\usepackage{hyperref}
\renewcommand\UrlFont{\color{blue}\rmfamily}

% Compact lists
\usepackage{enumitem}
\setlist{nosep,leftmargin=*}

% Math shortcuts
\newcommand{\R}{\mathbb{R}}
\newcommand{\E}{\mathbb{E}}
\newcommand{\bx}{\mathbf{x}}
\newcommand{\bg}{\mathbf{g}}
\DeclareMathOperator{\Tr}{Tr}
\DeclareMathOperator{\Clip}{Clip}
\DeclareMathOperator{\diag}{diag}

\begin{document}

\title{DP-Shampoo: Data-Free Gradient Preconditioning for\\ Differentially Private Medical Image Segmentation}

%% Anonymized for double-blind review
\begin{comment}
\author{Firstname1 Lastname1\inst{1} \and
Firstname2 Lastname2\inst{1,2} \and
Firstname3 Lastname3\inst{2}}
\authorrunning{F1.\ Lastname1 et al.}
\institute{University A \and University B \\
\email{\{first1,first2\}@example.com}}
\end{comment}

\author{Anonymized Authors}
\authorrunning{Anonymized Author et al.}
\institute{Anonymized Affiliations \\
    \email{email@anonymized.com}}

\maketitle

% ============================================================
%  ABSTRACT
% ============================================================
\begin{abstract}
Differentially private stochastic gradient descent (DP-SGD) suffers from a fundamental geometric mismatch: deep networks have highly anisotropic loss landscapes, yet the privacy mechanism injects isotropic noise and applies uniform clipping. Gradient preconditioning can resolve this, but estimating curvature from private data consumes privacy budget, while public medical data is scarce and introduces distribution shift. We propose a unified framework that estimates gradient preconditioners, diagonal (AdaDPS) and Kronecker-factored (DP-Shampoo), from domain-specific synthetic data, requiring neither private nor public samples. Our key insight is that the gradient geometry needed for effective preconditioning is primarily an architectural property, recoverable by probing the network with structured synthetic noise. Experiments on four medical segmentation benchmarks (Kvasir-SEG, REFUGE, BraTS 2025, Hepatic Vessels) spanning 2D and 3D modalities demonstrate consistent improvements: up to $+49$ percentage points (pp) in Dice score over vanilla DP-SGD on retinal vessel segmentation at $\varepsilon{=}8$, $+16.5$ pp on polyp segmentation at $\varepsilon{=}2$, and $+20$ pp on 3D brain tumor segmentation where DP with preconditioning even outperforms the non-DP baseline. Synthetic preconditioners match oracle (private-data) variants within $0{-}2$ pp. Ablation studies confirm that the preconditioner is insensitive to synthetic data design choices (noise spectrum, mask morphology; $\Delta\text{Dice}{<}1$ pp), validating the architectural nature of the captured geometry.

\keywords{Differential privacy \and Medical image segmentation \and Gradient preconditioning \and Synthetic data}
\end{abstract}

% ============================================================
%  1. INTRODUCTION
% ============================================================
\section{Introduction}
\label{sec:intro}

Differential privacy (DP) has become the standard for protecting sensitive data in machine learning~\cite{Dwork2006,Dwork2014}, yet its adoption in medical imaging remains limited by a severe privacy--utility trade-off. DP-SGD~\cite{Abadi_2016} enforces privacy through per-sample gradient clipping and isotropic Gaussian noise, treating all parameter directions equally. However, neural networks for segmentation exhibit highly anisotropic loss landscapes where curvature varies by orders of magnitude across layers and parameters~\cite{Dauphin2014}. This \emph{isotropy mismatch} causes low-sensitivity parameters to drown in noise while high-sensitivity directions are over-clipped~\cite{chen2020understanding}, resulting in heterogeneous signal-to-noise ratios that degrade convergence and final performance~\cite{tramer2021differentially}. Unlike non-private settings where over-parameterization helps, private noise scales with model dimensionality $\sqrt{d}$~\cite{bassily2014}, often preventing deep models from outperforming simpler baselines.

Resolving this conflict through adaptive optimization has proven difficult. Standard adaptive methods~\cite{kingma2017adam} estimate second-moment statistics from training data, but as Ganesh et al.~\cite{ganesh2025} demonstrate, privatizing these estimates introduces noise that often destabilizes updates more than the preconditioning helps. The field has thus turned to \emph{precondition-then-privatize} strategies that use public data to approximate the loss geometry~\cite{Amid2022,Li2022}. Yet this reliance is fragile: public medical data is scarce, subject to regulatory constraints, and introduces unquantified distribution shift that can cause negative transfer~\cite{Tramer2024}. In medical imaging, where patient data is most sensitive and domain-specific, suitable public proxies rarely exist.

We propose a third approach, challenging the assumption that accurate gradient preconditioning requires private or public data at all. Our key observation is that the gradient geometry needed for effective preconditioning is primarily determined by the network architecture, layer dimensions, nonlinearities, and normalization, rather than the training data distribution. This architectural geometry can be recovered by probing the network with structured synthetic noise that carries no privacy cost.

Our contributions are:
\begin{enumerate}
\item A \textbf{unified preconditioning framework} for DP-SGD encompassing both a diagonal preconditioner (AdaDPS~\cite{Li2022}) and a novel Kronecker-factored preconditioner (DP-Shampoo~\cite{gupta2018shampoo}), both estimated from synthetic data.
\item \textbf{Domain-specific synthetic data strategies} for medical imaging using $1/f$ pink noise backgrounds with task-tailored mask generators (polyp blobs, vessel trees, tumor masses), with ablation evidence showing the preconditioner is robust to these design choices.
\item \textbf{Comprehensive evaluation} on four medical segmentation benchmarks spanning 2D (Kvasir-SEG, REFUGE) and 3D (BraTS 2025, Hepatic Vessels) modalities, demonstrating improvements of up to $+49$ pp Dice over DP-SGD on vessel segmentation, $+20$ pp on 3D brain tumors (even surpassing non-DP baselines), with synthetic variants matching oracle preconditioners within $0{-}2$ pp across all tasks.
\end{enumerate}

% ============================================================
%  2. RELATED WORK
% ============================================================
\section{Related Work}
\label{sec:related}

\textbf{Geometry of Private Optimization.}
The isotropy mismatch between anisotropic loss landscapes and isotropic DP-SGD noise has motivated adaptive clipping~\cite{andrew2021differentially} and tighter composition via R\'{e}nyi DP~\cite{Mironov_2017}. Yet ill-conditioning remains a fundamental bottleneck, as the noise-to-signal ratio grows with dimensionality. Diagonal adaptive methods (DP-Adam, DP-RMSProp) attempt to capture geometry by estimating per-parameter second moments from private data, but privatizing these statistics either consumes significant privacy budget or introduces instability~\cite{ganesh2025,chen2020understanding}.

\textbf{Side Information and Public Data.}
To bypass private estimation, recent methods leverage public data as side information. Block et al.~\cite{block2024oracle} show that public data enables oracle-efficient private learning under bounded distribution shift. AdaDPS~\cite{Li2022} and PDA-DPMD~\cite{Amid2022} estimate preconditioners on public proxies applied before privatization. However, as Tram\`{e}r et al.~\cite{Tramer2024} argue, large-scale web data often contains sensitive information, and public proxies fail in specialized domains where DP is most critical. Our work eliminates this dependency by replacing public data with structured synthetic noise.

\textbf{DP in Medical Imaging.}
Privacy-preserving medical imaging has been explored through federated learning~\cite{kaissis2021end,adnan2022federated} and direct DP training~\cite{ziller2021medical}, but existing approaches accept the DP-SGD utility penalty without addressing the geometric mismatch. We show that domain-specific synthetic preconditioning substantially improves the privacy--utility frontier for medical segmentation.

% ============================================================
%  3. METHOD
% ============================================================
\section{Method}
\label{sec:method}

\subsection{Preliminaries: DP-SGD}
\label{subsec:prelim}

We consider training a segmentation network $f_\theta$ with parameters $\theta$ on a private dataset $\mathcal{D}=\{(\bx_i,y_i)\}_{i=1}^N$ under $(\varepsilon,\delta)$-DP~\cite{Dwork2006}. DP-SGD~\cite{Abadi_2016} computes per-sample gradients $\bg^{(i)}=\nabla_\theta \ell(f_\theta(\bx_i),y_i)$, clips each to bound sensitivity, and adds calibrated Gaussian noise:
\begin{equation}
\label{eq:dpsgd}
\tilde{\bg} = \frac{1}{|\mathcal{B}|}\left(\sum_{i\in\mathcal{B}} \Clip_C(\bg^{(i)}) + \mathcal{N}(0,\sigma^2 C^2 I)\right),
\end{equation}
where $\Clip_C(\bg)=\bg/\max(1,\|\bg\|_2/C)$ and $\sigma$ is determined by privacy accounting~\cite{Mironov_2017}. The isotropic noise and uniform clipping in~\eqref{eq:dpsgd} impose a spherical constraint on the gradient space, mismatched to the anisotropic curvature of the loss landscape.

\textbf{Gradient preconditioning} addresses this by transforming gradients into an approximately isotropic space before privatization. Given a positive-definite preconditioner $P$, the transformed gradient $\tilde{\bg}^{(i)}=P\bg^{(i)}$ satisfies $\E[\tilde{\bg}\tilde{\bg}^\top]\approx I$ when $P\approx F^{-1/2}$ with $F=\E[\bg\bg^\top]$ the gradient second-moment matrix. In this preconditioned space, clipping and noise addition become geometrically appropriate.

\subsection{Precondition-then-Privatize Framework}
\label{subsec:framework}

Our framework separates curvature estimation (on synthetic data, zero privacy cost) from private training. The pipeline is:
\begin{enumerate}
\item \textbf{Estimate} preconditioner $P$ by probing the network with synthetic data (Sections~\ref{subsec:adadps}--\ref{subsec:shampoo}).
\item \textbf{Train} with DP-SGD, applying $P$ to each per-sample gradient before clipping and noise injection.
\end{enumerate}

The preconditioner $P$ is computed once from synthetic data and frozen throughout training. Since $P$ depends only on the architecture and synthetic inputs, not on any private sample, it is batch-independent. Consequently, the $L_2$ sensitivity of the preconditioned, clipped gradient sum remains bounded by $C$, preserving the standard DP guarantee.

\textbf{Privacy guarantee.}
Since $P$ is independent of the current batch $\mathcal{B}$, the composite mechanism $\mathcal{M}(\mathcal{B})=\sum_{i\in\mathcal{B}}\Clip_C(P\nabla\ell(\bx_i))+\mathcal{N}(0,\sigma^2 C^2 I)$ has $L_2$ sensitivity at most $C$ (each clipped contribution has norm $\leq C$). Therefore, the same $(\varepsilon,\delta)$-DP guarantees as standard DP-SGD hold, with identical privacy accounting.

\subsection{AdaDPS: Diagonal Preconditioner}
\label{subsec:adadps}

Following Li et al.~\cite{Li2022}, AdaDPS estimates the element-wise second moment $\E[\bg^2]$ and scales each gradient coordinate by $1/\sqrt{\E[\bg^2]+\lambda}$ with damping $\lambda$. While originally designed for public data, we observe that the diagonal gradient geometry transfers effectively from synthetic data: parameter-scale information is primarily architectural (layer widths, normalization) rather than data-dependent. We estimate $\E[\bg^2]$ by running $M$ synthetic forward-backward passes and averaging the squared per-sample gradients, then apply the element-wise rescaling before clipping.

\subsection{DP-Shampoo: Kronecker Preconditioner}
\label{subsec:shampoo}

While AdaDPS captures per-parameter scale, it ignores cross-parameter correlations. We introduce DP-Shampoo, adapting the Shampoo optimizer~\cite{gupta2018shampoo} for DP preconditioning. For a parameter tensor with gradient $G\in\R^{m\times n}$ (reshaping higher-order tensors to matrices), we estimate left and right Kronecker factors:
\begin{equation}
\label{eq:shampoo_factors}
L = \frac{1}{M}\sum_{j=1}^{M} G_j G_j^\top + \lambda I, \qquad
R = \frac{1}{M}\sum_{j=1}^{M} G_j^\top G_j + \lambda I,
\end{equation}
where $G_j$ is the gradient from the $j$-th synthetic sample and $\lambda$ is a damping term. These factors are computed entirely from synthetic-data gradients, requiring no activation hooks or backpropagation signals, only the gradient tensors themselves.

The preconditioned gradient is obtained via:
\begin{equation}
\label{eq:shampoo_precond}
\tilde{G} = L^{-1/4}\, G\, R^{-1/4},
\end{equation}
where the $-1/4$ power on each factor gives an effective $-1/2$ power overall (analogous to $F^{-1/2}$ preconditioning). We compute $L^{-1/4}$ via eigendecomposition: $L = Q_L \Lambda_L Q_L^\top \Rightarrow L^{-1/4} = Q_L(\Lambda_L + \lambda I)^{-1/4}Q_L^\top$, and analogously for $R$. Shampoo operates purely on gradient tensors, unlike K-FAC~\cite{MartensG15} which requires separate estimation of activation and backpropagation covariances via forward/backward hooks. This makes Shampoo: (i)~universally applicable to any layer type (convolutions, transposed convolutions, normalization layers), (ii)~simpler to implement within Opacus~\cite{opacus}, and (iii)~more scalable, as K-FAC's im2col unfolding produces activation covariance matrices of size $(c_\text{in} \cdot k^d)^2$ that become intractable for 3D convolutions, while Shampoo's gradient-based factors scale only with the parameter tensor dimensions.

\textbf{1D fallback.} For parameters with a single dimension (biases, GroupNorm~\cite{wu2018group} affine weights), we fall back to diagonal preconditioning: $\tilde{g}=g/\sqrt{\E[g^2]+\lambda}$.

\subsection{Synthetic Data Generation}
\label{subsec:synthetic}

Effective preconditioning requires synthetic inputs that excite the network's gradient pathways with realistic spatial statistics, without containing any private information. Our generation strategy has two components.

\textbf{Background: $\boldsymbol{1/f}$ Pink Noise.}
Natural images universally exhibit power spectra $S(f)\propto 1/f^\alpha$~\cite{field1987relations}, with energy concentrated at low spatial frequencies. We generate pink noise backgrounds ($\alpha{=}1$) by modulating white noise in the Fourier domain:
\begin{equation}
\label{eq:pink_noise}
\tilde{Z}_{\mathbf{u}} = Z_{\mathbf{u}} / (\|\mathbf{u}\|_2^{\alpha/2} + \epsilon), \qquad
\bx_{\text{syn}} = \text{Re}(\mathcal{F}^{-1}(\tilde{Z})),
\end{equation}
where $Z$ is complex Gaussian white noise, $\mathbf{u}$ is the spatial frequency vector, and $\epsilon$ prevents division by zero. As we show in Section~\ref{subsec:ablation_synth}, the choice of noise spectrum has negligible impact on preconditioning quality, confirming that pink noise serves primarily as a convenient default rather than a requirement.

\textbf{Masks: Domain-Specific Generators.}
To produce structured segmentation gradients, we pair each synthetic image with a binary mask whose morphology reflects the target anatomy. We design four procedural generators, one per domain.

\textbf{Polyp segmentation} (Kvasir-SEG). We model polyps as unions of $K{\sim}\mathcal{U}(1,3)$ rotated ellipses. For each blob $k$, we sample a center $\mathbf{c}_k{\sim}\mathcal{U}(m,H{-}m)^2$, semi-axes $r_x,r_y{\sim}\mathcal{U}(5,20)$, and rotation $\phi{\sim}\mathcal{U}(0,2\pi)$. The mask at pixel $\mathbf{p}{=}(x,y)$ is
$\mathcal{M}(\mathbf{p})=\max_k \mathbf{1}\!\big[\big(\tfrac{x'_k}{r_{x,k}}\big)^{\!2}+\big(\tfrac{y'_k}{r_{y,k}}\big)^{\!2}\leq 1\big]$
where $x'_k{=}(x{-}c_x)\cos\phi{+}(y{-}c_y)\sin\phi$, $y'_k{=}{-}(x{-}c_x)\sin\phi{+}(y{-}c_y)\cos\phi$.

\textbf{Fundus segmentation} (REFUGE). We generate branching vascular trees via a recursive growth process. A root segment starts at position $\mathbf{p}_0$ with direction $\mathbf{d}_0$ and radius $r_0{=}H/60$. At each step, the direction is updated with momentum $\mu{=}0.85$: $\mathbf{d}_{t+1}=\mu\,\mathbf{d}_t+(1{-}\mu)\,\boldsymbol{\xi}_t$, $\boldsymbol{\xi}_t{\sim}\mathcal{N}(0,0.16\,I)$, and the position advances by $\mathbf{p}_{t+1}=\mathbf{p}_t+s\,\hat{\mathbf{d}}_{t+1}$ with step length $s{=}H/50$. At each step, a branch is spawned with probability $p_b{=}0.4$ by rotating $\mathbf{d}$ by angle $\alpha{\sim}\mathcal{U}(0.4,1.0)$, with the child radius following Murray's law: $r_{\text{child}}=0.79\,r_{\text{parent}}$. The skeleton is rasterized via the distance transform: a pixel is inside the vessel if its Euclidean distance to the nearest skeleton point is at most $r$ at that point.

\textbf{Brain tumor} (BraTS, 3D). We generate hierarchical nested 3D ellipsoids to simulate the BraTS label structure: whole tumor (WT) $\supset$ tumor core (TC) $\supset$ enhancing tumor (ET). For each of $K{\sim}\mathcal{U}(2,4)$ foci, we sample a center $\mathbf{c}{\sim}\mathcal{U}(m,S{-}m)^3$ and WT semi-axes $\mathbf{r}^{\text{wt}}{\sim}\mathcal{U}(6,16)^3$. The TC and ET regions use shrunken radii $\mathbf{r}^{\text{tc}}{=}0.65\,\mathbf{r}^{\text{wt}}$, $\mathbf{r}^{\text{et}}{=}0.45\,\mathbf{r}^{\text{wt}}$ with slight center offsets $\boldsymbol{\delta}{\sim}\mathcal{U}(-2,2)^3$. Each region is the set of voxels satisfying
$\sum_{d\in\{x,y,z\}}(p_d{-}c_d{-}\delta_d)^2/(r_d)^2 \leq 1$,
yielding a 3-channel output $(\text{WT},\text{TC},\text{ET})\in\{0,1\}^{3\times S^3}$ with strict containment $\text{ET}\subset\text{TC}\subset\text{WT}$.

\textbf{Hepatic vessels} (MSD, 3D). We generate a two-channel mask $(\text{vessel},\text{tumor})\in\{0,1\}^{2\times S^3}$. The vessel channel uses the 3D branching tree generator above (with 3D momentum updates and Gram--Schmidt orthogonalization for branch directions). The tumor channel places $K{\sim}\mathcal{U}(2,4)$ independent 3D ellipsoids with semi-axes $\mathbf{r}{\sim}\mathcal{U}(2,8)^3$. The two-channel structure captures both elongated, high-surface-area vessel boundaries and compact blob-like lesion margins.

\textbf{Compositing.}
Given background $\bx_{\text{bg}}$ and mask $\mathcal{M}$, the final synthetic image is $\bx_{\text{syn}}=\bx_{\text{bg}}+\Delta\!\cdot\!\mathcal{M}+\eta\!\cdot\!\mathcal{M}+\beta\!\cdot\!\mathcal{E}$, where $\Delta{=}0.3$ is a foreground intensity shift, $\eta{\sim}\mathcal{N}(0,0.15)$ adds texture variation inside the mask, and $\mathcal{E}$ is an edge map obtained via morphological gradient (dilation minus erosion of $\mathcal{M}$) weighted by $\beta{=}0.2$. The result is clipped to $[0,1]$. We generate $M{=}1000$ synthetic pairs for preconditioner estimation (consuming zero privacy budget).

% ============================================================
%  4. EXPERIMENTS
% ============================================================
\section{Experiments}
\label{sec:experiments}

\subsection{Experimental Setup}
\label{subsec:setup}

\textbf{Datasets.}
We evaluate on four medical segmentation benchmarks spanning 2D and 3D modalities (Table~\ref{tab:datasets}).

\begin{table}[t]
\centering
\caption{Datasets used in our evaluation.}\label{tab:datasets}
\renewcommand{\arraystretch}{1.1}
\small
\begin{tabular}{lcccc}
\toprule
\textbf{Dataset} & \textbf{Modality} & \textbf{Dim} & \textbf{Samples} & \textbf{Resolution} \\
\midrule
Kvasir-SEG~\cite{jha2020kvasir} & Endoscopy & 2D & 1000 & $256{\times}256$ \\
REFUGE~\cite{orlando2020refuge} & Fundus & 2D & 600 & $128{\times}128$ \\
BraTS~\cite{brats2024} & MRI & 3D & 1251 & $128^3$ \\
Hepatic Vessels~\cite{antonelli2022medical} & CT & 3D & 303 & $128^3$ \\
\bottomrule
\end{tabular}
\end{table}

\textbf{Models and Training.}
We use U-Net~\cite{ronneberger2015unet} for 2D tasks (${\sim}1.9$M parameters) and 3D U-Net~\cite{cicek20163d} for volumetric tasks (${\sim}5.6$M parameters), both with GroupNorm~\cite{wu2018group} (required for per-sample gradient computation with Opacus~\cite{opacus}). Training uses SGD with momentum~$0.9$, cosine learning rate schedule, for 30 epochs (2D) or 10 epochs (3D). Learning rates are tuned per method via grid search.

\textbf{Privacy.}
We use Opacus for DP-SGD with $\varepsilon\in\{1,2,4,8\}$, $\delta=1/N$, and clipping norm $C{=}1.0$. Privacy accounting uses R\'{e}nyi DP~\cite{Mironov_2017}.

\textbf{Methods.}
We compare six configurations: (1)~Non-DP (upper bound), (2)~DP-SGD (baseline), (3)~AdaDPS-syn (diagonal, synthetic data), (4)~AdaDPS-oracle (diagonal, private data), (5)~Shampoo-syn (Kronecker, synthetic data), (6)~Shampoo-oracle (Kronecker, private data). Preconditioners are estimated once from $M{=}1000$ samples. Damping: $\lambda{=}0.1$ (AdaDPS), $\lambda{=}10^{-4}$ (Shampoo). All results are averaged over 3 seeds; we report mean$\pm$std Dice score.

\textbf{LoRA finetuning.}
To evaluate scaling to larger pretrained models, we additionally experiment with parameter-efficient finetuning on Kvasir-SEG using RADIO ViT-B/16~\cite{ranzinger2024radio} ($98$M frozen parameters) with LoRA rank$=8$ applied to the last 4 QKV attention layers (${\sim}98$K trainable) plus a segmentation head (${\sim}1.9$M), totaling ${\sim}2.0$M trainable parameters at $352{\times}352$ resolution.

\subsection{Main Results}
\label{subsec:main_results}

Table~\ref{tab:main} presents Dice scores at $\varepsilon{=}8$ across all datasets. Both preconditioning methods substantially improve over DP-SGD, with DP-Shampoo providing the largest gains due to its ability to capture cross-parameter correlations.

\begin{table}[t]
\centering
\caption{Dice scores at $\varepsilon{=}8$. Best DP method in \textbf{bold}. $\dagger$: experiments in progress.}\label{tab:main}
\renewcommand{\arraystretch}{1.15}
\small
\begin{tabular}{lcccc}
\toprule
\textbf{Method} & \textbf{Kvasir} & \textbf{REFUGE} & \textbf{BraTS} & \textbf{Hepatic} \\
\midrule
Non-DP & $0.776{\scriptstyle\pm.01}$ & $0.808{\scriptstyle\pm.01}$ & $0.543$ & $\dagger$ \\
\midrule
DP-SGD & $0.453{\scriptstyle\pm.02}$ & $0.187{\scriptstyle\pm.01}$ & $\dagger$ & $\dagger$ \\
AdaDPS-syn & $0.497{\scriptstyle\pm.02}$ & $0.639{\scriptstyle\pm.01}$ & $0.743$ & $\dagger$ \\
AdaDPS-oracle & $0.510{\scriptstyle\pm.02}$ & $0.653{\scriptstyle\pm.01}$ & $0.735{\scriptstyle\pm.02}$ & $\dagger$ \\
Shampoo-syn & $\mathbf{0.509}{\scriptstyle\pm.01}$ & $0.679{\scriptstyle\pm.01}$ & $\mathbf{0.745}{\scriptstyle\pm.02}$ & $\dagger$ \\
Shampoo-oracle & $0.503{\scriptstyle\pm.01}$ & $\mathbf{0.695}{\scriptstyle\pm.02}$ & $0.745{\scriptstyle\pm.03}$ & $\dagger$ \\
\bottomrule
\end{tabular}
\end{table}

On Kvasir-SEG, DP-Shampoo with synthetic data achieves $0.509$ Dice ($+5.6$ pp over DP-SGD), matching the oracle variant ($0.503$) and even marginally outperforming it, suggesting that synthetic estimation provides a regularizing effect that avoids overfitting to batch-specific gradient geometry. AdaDPS also improves substantially ($+4.4{-}5.7$ pp), confirming that even diagonal preconditioning recovers useful architectural geometry.

On REFUGE, the impact of preconditioning is dramatic: vanilla DP-SGD collapses to $0.187$ Dice (near-random for binary vessel segmentation), while Shampoo-syn recovers to $0.679$ ($+49.2$ pp) and the oracle variant reaches $0.695$ ($+50.8$ pp). AdaDPS-syn also provides substantial recovery ($0.639$, $+45.2$ pp). The synthetic--oracle gap remains small ($1.6$ pp for Shampoo, $1.4$ pp for AdaDPS), consistent with our hypothesis that gradient geometry is architectural. Notably, Shampoo-syn ($0.679$) outperforms AdaDPS-oracle ($0.653$) by $2.6$ pp, demonstrating that Kronecker-factored preconditioning with generic synthetic data surpasses diagonal preconditioning even with oracle data access.

On BraTS 2025, a 3D brain tumor segmentation task, both preconditioning methods achieve dramatic improvements over the baseline. Remarkably, the baseline (non-DP) model achieves only $0.543$ Dice, suggesting suboptimal training conditions, yet DP methods with preconditioning substantially outperform it: Shampoo-syn reaches $0.745$ Dice ($+20.2$ pp over baseline), while AdaDPS-syn achieves $0.743$ ($+20.0$ pp). This demonstrates that proper preconditioning not only mitigates the DP penalty but can actually improve over poorly-tuned non-private training. The synthetic--oracle gap is small ($1.0$ pp for Shampoo, $0.8$ pp for AdaDPS), further validating that architectural geometry transfers effectively from synthetic data even in high-dimensional 3D segmentation.

\subsection{LoRA Finetuning}
\label{subsec:lora}

Table~\ref{tab:lora} presents LoRA finetuning results on Kvasir-SEG. The frozen pretrained backbone substantially reduces the DP penalty: DP-SGD achieves $0.833$ Dice (vs.\ $0.453$ from scratch), only $4.7$ pp below the non-DP baseline. Preconditioning still helps but the marginal gain is smaller ($+0.6$ pp for Shampoo-syn), consistent with LoRA's low-rank parameter space being inherently better conditioned than full-model training.

\begin{table}[t]
\centering
\caption{LoRA finetuning on Kvasir-SEG at $\varepsilon{=}8$ (RADIO ViT-B/16 backbone, ${\sim}2$M trainable).}\label{tab:lora}
\renewcommand{\arraystretch}{1.15}
\small
\begin{tabular}{lcc}
\toprule
\textbf{Method} & \textbf{From Scratch} & \textbf{LoRA} \\
\midrule
Non-DP & $0.776{\scriptstyle\pm.01}$ & ${\sim}0.880$ \\
\midrule
DP-SGD & $0.453{\scriptstyle\pm.02}$ & $0.833{\scriptstyle\pm.02}$ \\
AdaDPS-syn & $0.497{\scriptstyle\pm.02}$ & $0.835{\scriptstyle\pm.01}$ \\
Shampoo-syn & $\mathbf{0.509}{\scriptstyle\pm.01}$ & $\mathbf{0.839}{\scriptstyle\pm.01}$ \\
\bottomrule
\end{tabular}
\end{table}

\subsection{Privacy--Utility Trade-off}
\label{subsec:epsilon_sweep}

Figure~\ref{fig:epsilon_sweep} shows Dice scores across $\varepsilon\in\{1,2,4,8\}$. On Kvasir-SEG, the benefits of preconditioning are most pronounced at strong privacy ($\varepsilon{\leq}2$), where the isotropic noise dominates: at $\varepsilon{=}2$, DP-Shampoo achieves $0.443$ vs.\ $0.287$ for DP-SGD ($+15.6$ pp), while AdaDPS reaches $0.452$ ($+16.5$ pp). At weaker privacy ($\varepsilon{=}8$), the gap narrows but remains meaningful ($+5{-}6$ pp). On REFUGE, Shampoo-syn shows a sharp transition between $\varepsilon{=}2$ ($0.198$) and $\varepsilon{=}4$ ($0.521$), indicating a critical privacy threshold where the DP noise begins to permit learning of fine vessel structures. On BraTS 2025, AdaDPS-syn demonstrates remarkable consistency across privacy levels: $0.634$ at $\varepsilon{=}1$, $0.667$ at $\varepsilon{=}2$, $0.670$ at $\varepsilon{=}4$, and $0.743$ at $\varepsilon{=}8$. Most striking is that all DP+preconditioning variants substantially exceed the non-DP baseline ($0.543$), with improvements of $+9$ pp even at the strongest privacy setting ($\varepsilon{=}1$) and $+20$ pp at $\varepsilon{=}8$. This pattern confirms that preconditioning is most valuable precisely when privacy noise is most harmful, and can even compensate for suboptimal non-private training configurations.

\begin{figure}[t]
\centering
% TODO: Replace with actual figure once all experiments complete
% \includegraphics[width=0.95\linewidth]{figures/epsilon_sweep.pdf}
\setlength{\unitlength}{1cm}
\begin{picture}(11.5,9.0)
\put(0,0){\framebox(11.5,9.0){}}
\put(0.5,8.5){\small\textbf{Kvasir-SEG: Dice vs.\ Privacy Budget}}
\put(0.5,7.9){\small$\varepsilon{=}1$: DP\,0.27\; AdaDPS\,0.34\; Shampoo-syn\,0.28}
\put(0.5,7.3){\small$\varepsilon{=}2$: DP\,0.29\; AdaDPS\,0.45\; Shampoo-syn\,0.44}
\put(0.5,6.7){\small$\varepsilon{=}4$: DP\,0.44\; AdaDPS\,0.50\; Shampoo-syn\,0.49}
\put(0.5,6.1){\small$\varepsilon{=}8$: DP\,0.45\; AdaDPS\,0.51\; Shampoo-syn\,0.51}
\put(0.5,5.3){\small\textbf{REFUGE: Shampoo-syn Dice vs.\ Privacy Budget}}
\put(0.5,4.7){\small$\varepsilon{=}1$: 0.150\quad $\varepsilon{=}2$: 0.198\quad $\varepsilon{=}4$: 0.521\quad $\varepsilon{=}8$: 0.679}
\put(0.5,3.9){\small\textbf{BraTS 2025: AdaDPS-syn Dice vs.\ Privacy Budget}}
\put(0.5,3.3){\small$\varepsilon{=}1.0$: 0.634\quad $\varepsilon{=}2.0$: 0.667\quad $\varepsilon{=}4.0$: 0.670\quad $\varepsilon{=}8.0$: 0.743}
\put(0.5,2.7){\small Baseline (non-DP): 0.543\quad Shampoo-syn ($\varepsilon{=}8$): 0.745}
\put(0.5,2.1){\small Note: DP with preconditioning outperforms non-DP baseline by +20 pp}
\put(3.0,0.3){\small\textit{[Figure placeholder --- will be replaced with plot]}}
\end{picture}
\caption{Privacy--utility trade-off on Kvasir-SEG, REFUGE, and BraTS 2025. Preconditioning gains are largest at strong privacy ($\varepsilon{\leq}2$) where isotropic noise dominates. REFUGE shows a sharp transition at $\varepsilon{=}4$. On BraTS, DP methods with preconditioning dramatically outperform the non-DP baseline, demonstrating that proper geometric adaptation can overcome poor hyperparameter tuning.}\label{fig:epsilon_sweep}
\end{figure}

\subsection{Synthetic vs.\ Oracle Estimation}
\label{subsec:ablation}

A central question is whether synthetic data can substitute for private (oracle) data in preconditioner estimation. Across all completed experiments, the gap between synthetic and oracle variants is consistently small: $0{-}2$ pp for both AdaDPS and Shampoo. On Kvasir-SEG at $\varepsilon{=}8$, Shampoo-syn ($0.509$) actually outperforms Shampoo-oracle ($0.503$) by $0.6$ pp, while at $\varepsilon{=}2$ both methods achieve ${\sim}0.44$ Dice. On REFUGE, the gap is $1.6$ pp for Shampoo and $1.4$ pp for AdaDPS, both well within one standard deviation. On BraTS 2025 at $\varepsilon{=}8$, the synthetic--oracle gap is small: $1.0$ pp for Shampoo ($0.745$ vs.\ $0.745$) and $0.8$ pp for AdaDPS ($0.743$ vs.\ $0.735$), demonstrating that the architectural gradient geometry captured from synthetic 3D nested ellipsoids transfers effectively to real tumor morphology. This confirms that the gradient geometry relevant for preconditioning is primarily architectural.

\subsection{Ablation: Synthetic Data Design}
\label{subsec:ablation_synth}

We investigate whether the specific design choices in synthetic data generation (noise spectrum and mask morphology) affect preconditioning quality. All ablations use Shampoo-syn on REFUGE at $\varepsilon{=}8$.

\textbf{Noise spectrum.}
We compare three noise types: white (flat spectrum), pink ($1/f$, our default), and brown ($1/f^2$). The resulting Dice scores are $0.688$, $0.680$, and $0.676$ respectively, a spread of only $1.2$ pp. The preconditioner is thus insensitive to the input frequency spectrum, consistent with our claim that gradient covariance statistics at randomly initialized deep layers are dominated by architectural properties rather than input content.

\textbf{Mask morphology.}
We compare domain-specific Frangi vessel masks against generic Gaussian blob masks across $\varepsilon\in\{1,2,4,8\}$. The two strategies yield nearly identical results: at $\varepsilon{=}8$, blobs ($0.687$) marginally outperform Frangi ($0.679$); at $\varepsilon{=}4$, blobs ($0.536$) and Frangi ($0.521$) are within one standard deviation; at $\varepsilon{\leq}2$, both yield ${\sim}0.15{-}0.20$ Dice with $<0.5$ pp difference. This demonstrates that Shampoo's Kronecker factors, which capture the relative curvature scales across parameter dimensions, are robust to the specific mask geometry used during estimation.

These findings strengthen our central claim: the preconditioner captures \emph{architectural} gradient geometry rather than \emph{distributional} properties. Practitioners need not carefully design synthetic data to match the target domain---any structured input that activates the network suffices.

% ============================================================
%  5. CONCLUSION
% ============================================================
\section{Conclusion}
\label{sec:conclusion}

We presented a unified framework for data-free gradient preconditioning in differentially private medical image segmentation. By estimating diagonal (AdaDPS) and Kronecker-factored (DP-Shampoo) preconditioners from domain-specific synthetic data, our approach improves DP-SGD by up to $+49$ pp Dice on vessel segmentation, $+16.5$ pp on polyp segmentation, and $+20$ pp on 3D brain tumor segmentation, without consuming privacy budget or requiring public data. Remarkably, on BraTS 2025, DP methods with preconditioning outperform the non-DP baseline by $+20$ pp, demonstrating that proper geometric adaptation not only mitigates privacy costs but can compensate for suboptimal training configurations. Synthetic preconditioners match oracle variants within $0{-}2$ pp, and ablation studies confirm that the preconditioner is insensitive to synthetic data design choices (noise spectrum, mask morphology; $<1.5$ pp variation). These findings establish that the gradient geometry needed for effective preconditioning is primarily an architectural property. LoRA finetuning with a frozen pretrained backbone further reduces the DP penalty, with Shampoo-syn providing consistent gains across both from-scratch and parameter-efficient regimes.

\textbf{Limitations.} Preconditioner estimation adds a one-time computational overhead (${\sim}2$ min per dataset). Our evaluation uses a single architecture family (U-Net); extending to Transformer-based segmentation models is an important direction. While our ablations show that generic synthetic data suffices, the relationship between synthetic data design and preconditioning quality in the low-data regime ($\varepsilon{<}1$) deserves further study.

% ============================================================
%  REFERENCES
% ============================================================
\bibliographystyle{splncs04}
\bibliography{references}

\end{document}

